%%% File:       a02_Preface.tex
%%% Component:  Preface
%%% Author:     Joaquín Ataz-López et al
%%% Begun:      March 2020
%%% Concluded:  August 2020
%%%
%%% Edited with: Emacs + AUCTeX - and at times with vim + context-plugin
%%%

\startcomponent a02_Preface.tex

\switchtobodyfont[sl]

% 本节的脚注将使用符号而非数字编号：
\setupnotation
  [footnote]
  [numberconversion=set 2]

\starttitle
  [title=前言]

\writetolist[chapter]{}{前言}

\setupheadertexts
  [\tfx 前言]
  [\tfx{\romannumerals{\pagenumber}}]

{\it 一本不算太短的 \ConTeXt\ Mark~IV 入门}（有时也被称为 ANSS，至少在英语 \ConTeXt\ 用户中如此）由 Joaquín Ataz-López 所著，是备受 \ConTeXt\ 用户推崇的参考文档，对新用户尤其有价值，无论他们之前是否有过 plain \TeX、\LaTeX\ 的经验，还是对 \TeX\ 家族完全陌生。该书已被从西班牙语原版翻译成其他五种语言，足见其广泛吸引力。

时代在前进，\ConTeXt\ 也在发展。我敢大胆推测，到 2026 年，大多数 \ConTeXt\ 用户都在使用 \ConTeXt~LMTX\null。尽管 Mark~IV 和 LMTX 版本的 \ConTeXt\ 非常相似，但 ANSS 中的一些信息已不能反映 \ConTeXt\ LMTX\null 的当前行为。此外，LMTX 中还有值得提及的新功能。

Joaquín 慷慨地支持为其著作创建新 \quotation{社区版}的想法。同时，对标题和署名的（微小）调整是合适的，以反映目标 \ConTeXt\ 版本以及更广泛的作者群体（尽管 Joaquín 的工作始终是这本新书的基础）。

{\bf 本}书的 \ConTeXt\ 源代码当前位于 \from[SourceCode]。欢迎所有 \ConTeXt\ 用户贡献。对于熟悉 {\tt git\/} 的用户，欢迎提交拉取请求。在 Codeberg 网站上开启一个 \quotation{议题}是另一种贡献方式。第三种选择是将补丁通过电子邮件发送给其中一位维护者。最后，对本书的重大更改（例如关于轻微或重大重组的想法）最好通过 \ConTeXt\ 邮件列表处理。

现有书籍的改编值得一篇新的前言。我决定这样做，以推动项目前进。然而，我必须明确无误地强调，我在贡献者中并无特殊地位；我不是 \ConTeXt\ 专家，也没有撰写如此规模技术文档的经验。我希望 \ConTeXt\ 社区的许多成员能根据自己的技能和空余时间做出贡献，为所有 \ConTeXt\ 用户提供一份有用且最新的文档。

ANSS 的前言被收录到本书中，原因有几个，这里我只提四个。
首先，它解释了 ANSS 开发背后的原理；\Conjecture\ 这些原理依然存在，并部分解释了为什么这本 \quotation{ANSS-CE} 书被认为有价值。其次，它解释了左边距中看到的两幅图像的含义，读者肯定想知道这些信息。第三，那篇前言包含了一些关于 \TeX\ 和 \ConTeXt\ 的历史背景 \Doubt，我相信值得花几分钟阅读。第四，解释了 ANSS（以及 ANSS-CE，除非未来某个时候进行重大更改）的整体结构，这有助于读者更好地理解本书试图做什么，以及它试图如何做。

最后，除本前言外，所有使用“我”的地方均指 Joaquín Ataz-López。随着本书的发展，“我”的数量可能会减少，但本书的起源将始终与 ANSS\null 同在。

\blank

\rightaligned{Jim Diamond}\\
\rightaligned{2026年3月}

\page

\stoptitle


\starttitle
  [title=贡献者列表]
\writetolist[chapter]{}{贡献者列表}

本节列出已知为本书做出贡献（无论大小）的人员；无疑遗漏了许多人。如果您的名字应在此列，请按照前言中的说明报告遗漏。（如果您的名字在此，而您谦虚地希望将其改为 \quotation{匿名 \#<n>}，请告知维护者。）

\startlines
{\bf Joaquín Ataz-López}（ANSS 作者）
{\bf 匿名 \#1}（西班牙语原版书的英文译者）

Aditya Mahajan
Florent Michel
Jim Diamond
Hans Hagen
Henning Hraban Ramm
Max Chernoff
Pablo Rodriguez
Wolfgang Schuster
\stoplines

\stoptitle

\starttitle
  [title=《一本不算太短的 \ConTeXt\ Mark~IV 入门》前言\footnote{本前言起初打算翻译/改编自 \quotation{《\TeX\ 书》} 的前言，该文档解释了 {\em 关于 \TeX\ 你需要知道的一切}。最终，我不得不偏离这一初衷；然而，我保留了一些片段，希望对于那些了解它的人，能唤起一些 {\em 回响}。}]

\writetolist[chapter]{}{《一本不算太短的 \ConTeXt\ Mark~IV 入门》前言}

\setupnotation
  [footnote]
  [numberconversion=n]

\setupheadertexts
  [\tfx ANSS 前言]
  [\tfx{\romannumerals{\pagenumber}}]

亲爱的读者，这是一份关于 \ConTeXt 的文档，\ConTeXt\ 是一个源自 \TeX\ 的排版系统，而 \TeX\ 又是另一个由 {\sc Donald E.~Knuth} 于 1977 年至 1982 年间在斯坦福大学创建的排版系统。

\ConTeXt\ 旨在创建具有极高排版质量的文档——无论是纸质文档还是设计用于计算设备屏幕上显示的文档。它不是一个文字处理器或文本编辑器，而是如前所述，一个 {\em 系统}，或者换句话说，一套旨在排版文档的工具，排版被理解为文档不同元素在页面或屏幕上的图形布局和可视化。总之，\ConTeXt\ 旨在提供所有必要的工具，以使文档具有最佳的外观。其理念是能够生成不仅写得好，而且 \quotation{优美}的文档。在这方面，我们可以引用 {\sc Donald E.~Knuth} 在介绍 \TeX\ （\ConTeXt\ 所基于的系统）时所写的话：

\startnarrower\it

  如果你只是想制作一份勉强过得去的文档——可以接受、基本可读但并不真正优美——一个更简单的系统通常就足够了。使用 \TeX\ 的目标是制作出 {\em 最优}的质量；这需要更注意细节，但你不会觉得多花点功夫有多难，而且你会对成品感到特别自豪。

\stopnarrower

当我们用 \ConTeXt\ 准备稿件时，我们会精确指示如何将其转换为页面（或屏幕），其排版质量可与世界上最好的印刷厂相媲美。要做到这一点，一旦我们学会了这个系统，我们所需的工作量只比在任何文字处理器或文本编辑器中正常录入文档稍微多一点。事实上，一旦我们在处理 \ConTeXt\ 上获得了一定的熟练度，考虑到文档的主要格式细节是在 \ConTeXt\ 中全局描述的，并且我们处理的是文本文件——一旦习惯了，这是一种更自然的创建和编辑文档的方式，我们的总工作量可能会更少；此外，这类文件比文字处理器那种庞大的二进制文件要轻得多，也更容易处理。

关于 \ConTeXt\ 有相当多的文档，几乎全是英文。我们可能认为是 \ConTeXt\ {\em 官方}发行版的东西——称为 \suite-{}\footnote{在本文第一版起草时，那里的陈述是符合事实的；但在 2020 年春天，ConTeXt wiki 进行了更新，从那时起，我们必须认为 \ConTeXt\ 的 \quotation{官方}发行版已经变成了 LMTX\null。然而，对于第一次接触 ConTeXt 世界的人来说，我仍然建议使用 \suite-{}，因为它是一个更稳定的发行版。
  \in{附录}[installation_suite] 解释了如何安装任一发行版。}——例如，包含大约 180 个 PDF 文档（大部分是英文，但也有一些荷兰语和德语），包括手册、示例和技术文章；在 Pragma ADE 网站（孕育了 \ConTeXt\ 的公司）上，（在我 2020 年 5 月统计的那天）有 224 个可下载的文档，其中大部分随 \suite- 一起分发，但也有一些是其他的。尽管如此，这些庞大的文档对于学习 \ConTeXt\ 并不是特别有用，因为通常这些文档的目标读者并不是对系统一无所知但想学习它的人。在 \suite- 称为 \quotation{手册}的 56 个 PDF 文件中，只有一个假设读者对 \ConTeXt\ 一无所知。这份文档名为 \quotation{\ConTeXt\ Mark~IV，一瞥}。然而，这份文档，正如其名称所示，仅限于介绍系统并解释如何使用 \ConTeXt\ 做某些事情。如果后续有一本结构更清晰、更系统的参考手册，它会是一个很好的入门介绍。这本参考手册并不存在，而且与 \ConTeXt\ \quotation{一瞥}相关的文档与其他文档之间的差距太大了。

2001 年，有一本参考手册被编写出来，可以在 \goto{Pragma ADE 网站}[url(pragma)] 上找到；但尽管有这个名字，一方面它并非旨在成为一本 {\em 完整的手册}，另一方面它针对的是 \ConTeXt\ 的旧版本（称为 Mark~II），因此已经相当过时。

2013 年，该手册被部分更新，但许多章节没有被重写，它包含的信息同时涉及 \ConTeXt\ Mark~II 和 \ConTeXt\ Mark~IV（当前版本），但并非总能完全清楚哪些信息指的是哪个版本。也许这就是为什么这本手册没有出现在 \suite- 包含的文档中。然而，尽管存在这些缺陷，在我们阅读了介绍性的 \quotation{\ConTeXt\ Mark~IV，一瞥}之后，该手册仍然是开始学习 \ConTeXt\ 的最佳文档。对于 \ConTeXt\ 入门同样非常有用的是其 \goto{wiki}[url(wiki)] 上的信息，在撰写本文时，该 wiki 正在重新设计，结构更加清晰，尽管它也将仅适用于 Mark~II 的解释与其他适用于 Mark~IV 或两个版本的解释混合在一起。这种缺乏区分的情况也出现在 \ConTeXt\ 命令的官方列表中\footnote{列表见 \in{节}[sec:qrc-setup-en]}，该列表没有指明哪些命令仅适用于两个版本之一。

基本上，本入门书是根据这里列出的四个信息来源编写的：\ConTeXt\ \quotation{一瞥}、2013 年手册、wiki 的内容以及命令的官方列表（其中包含每个命令允许的配置选项）；此外，当然还有我自己的测试和结论。所以，事实上这本入门书是研究工作的结果，有一段时间我倾向于将其命名为 \quotation{我所知道的 \ConTeXt\ Mark~IV} 或 \quotation{我学到的关于 \ConTeXt\ Mark~IV 的知识}。最终，我放弃了这些标题，因为尽管它们可能很真实，但我觉得它们可能会阻止某些人进入 \ConTeXt；可以肯定的是，尽管文档（在我看来）存在一些缺陷，但这确实是一个有用且多功能的工具，学习它所付出的努力无疑是值得的。通过使用 \ConTeXt，我们可以操作和配置文本文档，以实现不了解该系统的人根本无法想象的事情。

\startSmallPrint

  由于我的身份，我无法避免在整个文档中不时出现我对信息缺乏的抱怨。我不希望这被误解：我非常感激 \ConTeXt\ 的创造者们设计了如此强大的工具并将其提供给公众。只是我忍不住认为，如果它的文档得到改进，这个工具会受欢迎得多：人们需要投入大量时间来学习它，这与其内在难度（它确实有难度，但并不比其他类似工具更大——事实上反而更小）关系不大，而是由于缺乏清晰、完整和组织良好的信息来区分 \ConTeXt\ 的两个版本，解释每个版本中的功能，最重要的是，阐明每个命令、参数和选项的作用。

  诚然，这类信息需要大量的时间投入。但由于许多命令共享名称相似的选项，或许可以提供一种 {\em 选项词汇表}，这也有助于发现一些不一致之处，例如两个名称相同的选项做不同的事情，或者为了做同一件事，在不同的命令中使用不同名称的选项。

  至于第一次接触 \ConTeXt\ 的读者，请不要让我的抱怨打消您的积极性，因为尽管信息不足确实会增加学习所需的时间，但至少对于本入门书所涉及的内容，我已经投入了这些时间，读者无需再投入。仅凭从本入门书中学到的知识，读者将拥有一个工具，使他们能够以前所未有的轻松方式制作文档。

\stopSmallPrint

由于本文档中解释的内容很大程度上来自我自己的结论，因此很可能，尽管我已经亲自测试了所解释的大部分内容，但某些陈述或观点可能既不正确也不够正统。当然，我会感激读者能提供的任何修正、细微差别或澄清，这些可以发送到 \color[maincolor]{joaquin@ataz.org}。然而，为了减少我可能出错的情况，我尽量不涉及那些我没有找到任何信息且无法（或不想）亲自尝试的问题。有时这是因为我的测试结果不具结论性，有时则是因为我并非总是测试所有内容：\ConTeXt\ 拥有的命令和选项数量惊人，如果我必须测试所有内容，我将永远无法完成这本入门书。然而，有些情况下我无法避免 {\em 假设}某事，即做出我认为可能但并非完全确定的陈述。在这些情况下，段落左边距中会放置一个 \quote{推测} 图像。图像 \Conjecture{} 旨在图形化地表示假设。\footnote{我不是画的这幅图，而是从网上下载的
（\goto{https://es.dreamstime.com/}[url(https://es.dreamstime.com/icono-s\%C3\%B3lido-negro-para-la-conjetura-preocupaci\%C3\%B3n-y-duda-conjeturar-el-logotipo-image146773292)]），上面说它是免版税图像。} 在其他时候，我别无选择，只能承认我不知道某事，甚至对其没有合理的假设：在这种情况下，左边距中可见的图像\Doubt~旨在表示的不仅仅是推测或无知。\footnote{也是在互联网上找到的
（\goto{https://www.freepik.es/}[url(https://www.freepik.es/iconos-gratis/duda-cabeza_785036.htm)]），那里授权免费使用。} 但由于我向来不擅长图形表示，我不确定我选择的图像是否真的能传达这么多细微差别。

另一方面，这本入门书是从一个对 \TeX\ 和 \ConTeXt\ 都一无所知的读者的角度编写的，尽管我希望它也能对那些从 \TeX\ 或 \LaTeX\ （最流行的 \TeX\ 衍生品）转来、第一次接触 \ConTeXt\ 的人有所帮助。尽管如此，我意识到试图取悦如此多不同类型的读者，我有可能无法让任何人满意。因此，在不确定的情况下，我始终明确，本文档的主要受众是 \ConTeXt\ 的新手，刚刚进入这个迷人生态系统的新手。

成为 \ConTeXt\ 的新手并不意味着也是计算机工具使用的新手；尽管在本入门书中我不假设读者具有任何特定水平的计算机素养，但我确实假定了一定的 \quotation{合理素养}，这意味着，例如，对文字处理器和文本编辑器之间的区别有大致了解，知道如何创建、打开和操作文本文件，知道如何安装程序，知道如何打开终端并执行命令……仅此而已。

\startSmallPrint

  在我撰写这些文字时，重读本入门书的前面部分，我意识到有时我会跑题，深入到学习 \ConTeXt\ 所不需要的计算机问题中，这可能会吓跑新手，而其他时候我又忙于解释一些相当明显的事情，这可能会让有经验的读者感到无聊。我恳求两者的宽容。理性上，我知道一个完全初涉计算机化文本管理的人甚至几乎不知道 \ConTeXt\ 的存在，但从另一个角度来看，在我的专业环境中，我周围都是那些在使用文字处理器时不断与文本作斗争的人，他们做得相当好，但从未处理过文本文件本身，因此他们忽略了诸如文本文件使用什么编码或文字处理器和文本编辑器之间的区别等基本问题。

\stopSmallPrint

本手册是为对 \ConTeXt\ 或 \TeX\ 一无所知的人设计的，这意味着我包含了明显不是关于 \ConTeXt\ 而是关于 \TeX\ 的信息；但我认为没有必要用对读者不相关的信息来加重他们的负担，例如某个命令 {\em 实际上} 是有效的，但它到底是 \ConTeXt\ 命令还是属于 \TeX\ ；因此，只有在某些时候，当我认为有用时，我才会澄清某些命令实际上属于 \TeX\null 。

关于本文档的组织结构，材料分为三个部分：

\startitemize

  \item {\bf 第一部分}，包括前四章，提供了 \ConTeXt\ 的全局概述，解释它是什么以及我们如何使用它，展示了如何转换文档的第一个示例，以便随后解释 \ConTeXt\ 的一些基本概念以及与 \ConTeXt\ 源文件相关的某些问题。

  总的来说，这些章节面向的是迄今为止只知道如何使用文字处理器的读者。已经了解标记语言的读者可以跳过前面的这些章节；如果读者已经了解 \TeX\ 或 \LaTeX，他们也可以跳过第 3 章和第 4 章的大部分内容。尽管如此，我建议至少阅读：

  \startitemize

    \item 关于 \ConTeXt\ 命令的信息（第 3 章），特别是它的功能、配置方式，因为这是 \LaTeX\ 和 \ConTeXt\ 在构思和语法上的主要区别所在。由于本入门书仅涉及后者，这些区别并未明确指明，但是阅读本章的读者如果了解 \LaTeX\ 的工作原理，将立即理解这两种语言的语法差异，以及 \ConTeXt\ 允许我们配置和定制几乎所有命令工作方式的方法。

    \item 关于多文件 \ConTeXt\ 项目的信息（第 4 章），这与使用其他基于 \TeX\ 系统的工作方式不太相似。

  \stopitemize

  \item {\bf 第二部分}，包括第 5 章到第 9 章，重点关注我们认为的 \ConTeXt\ 文档的主要全局方面：

  \startitemize

    \item 主要影响文档外观的两个方面是页面尺寸和布局以及所用字体。第 5 章和第 6 章专门讨论这些问题。

    \startitemize

      \item 第一章聚焦于页面：尺寸、构成页面的元素、布局（即页面元素如何分布）等。出于系统性原因，与分页和允许我们影响分页的机制相关的更具体方面也在此讨论。

      \item 第 6 章解释与字体及其处理相关的命令。这里还包括对颜色使用和管理的基本解释，因为尽管颜色严格来说不是字体的 {\em 特征}，但它同样影响文档的外部外观。

    \stopitemize

    \item 第 7 章和第 8 章聚焦于文档的结构以及 \ConTeXt\ 为作者提供的编写结构良好的文档的工具。第 7 章侧重于严格意义上的结构（文档的结构划分），第 8 章侧重于这如何反映在目录中；不过，在解释这一点时，我们借此机会也解释了如何使用 \ConTeXt\ 生成各种索引，因为对于 \ConTeXt\ 来说，这些都归在 \quotation{列表}的概念下。

    \item 最后，第 9 章聚焦于引用，这是任何文档在需要引用文档其他部分（内部引用）或其他文档（外部引用）时的一个重要全局方面。对于后者，我们目前只对指向外部文档的引用（链接）感兴趣。这些 {\em 链接}（也可能出现在内部引用中）使我们的文档具有 {\em 交互性}，在本章中我们将解释 \ConTeXt\ 用于创建此类文档的一些特性。

  \stopitemize

  这些章节不需要按特定顺序阅读，除了第 8 章，如果先阅读第 7 章会更容易理解。无论如何，我已尽力确保当某个章节或小节中出现的问题在本文档的其他地方有讨论时，文本会包含提及并附上指向该讨论点的超链接。然而，我无法保证总是如此。

  \item 最后是 {\bf 第三部分}（第 10 章及以后），侧重于更详细的方面。它们不仅彼此独立，甚至其小节也是独立的（除了最后一章可能例外）。鉴于 \ConTeXt\ 包含的大量实用工具，这部分可能会非常广泛；但由于我的理解是，当读者到达这里时，他们已经准备好主动深入研究 \ConTeXt\ 的文档，所以我只包含了以下几章：

  \startitemize

    \item 第 10 章和第 11 章讨论我们可以称之为任何文本文档的 {\em 核心元素}：文本由字符组成，字符组成单词，单词排列成行，行又组成段落，段落之间由垂直间距分隔……显然，所有这些问题本可以放在一章中，但由于那样会太长，我将这个问题分成两章，一章讨论字符、单词和水平间距，另一章讨论行、段落和垂直间距。

    \item 第 12 章是一个 {\em 大杂烩}，处理文档中常见的元素和结构，主要是学术、科学或技术文档：脚注、结构化列表、描述、编号等。

    \item 最后，第 13 章聚焦于浮动对象，尤其是其中最具代表性的：插入文档的图像和表格。

  \stopitemize

  \item 入门书以三个 {\bf 附录} 结束。一个是关于安装 \ConTeXt，第二个附录包含几十个允许生成各种符号的命令——主要但不仅限于数学用途，第三个附录包含本文档中解释或提及的 \ConTeXt\ 命令的字母顺序列表。

\stopitemize

还有许多问题有待解释：处理引文和参考文献、撰写专业文本（数学、化学……）、与 XML 的连接、与 Lua 代码的接口、模式和基于模式的处理、使用 MetaPost 设计图形等等。这就是为什么，由于我没有包含对 \ConTeXt\ 的完整解释，也不打算这样做，我将本文档称为 \quotation{\ConTeXt\ Mark IV 入门}；并且我补充说这个入门书不算太短，因为显然情况如此：一份留下这么多未完成事项但已超过 300 页的文本绝不是简短的入门。这是因为我希望读者理解 \ConTeXt\ 的逻辑，或者至少是我所理解的逻辑。它并不声称是一本参考手册，而是一本自学指南，让读者准备好制作中等复杂度的文档（这包括大多数可能遇到的文档），并且最重要的是教会读者 {\em 想象} 可以用这个强大的工具做什么，并在现有文档中找到如何做的方法。本文档也不是 {\em 教程}。教程旨在逐步增加难度，以便逐步教授要学的内容；在这方面，我选择从第二部分开始，而不是按照难度顺序排列材料，以便更系统。但尽管它不是教程，我还是包含了非常多的示例。

对于某些读者来说，本文档的标题可能会让他们想起由 {\sc Oetiker, Partl, Hyna} 和 {\sc Schlegl} 编写的一份互联网上的文本，那是介绍 \LaTeX\ 世界的最佳文档之一。我指的是 \quotation{{\em 一本不算太短的 \LaTeX\ $2_{\epsilon}$ 入门}}。这并非巧合，而是一种致敬和感激：感谢那些撰写像那样的文本的人的慷慨工作，使得许多人能够开始使用像 \LaTeX\ 和 \ConTeXt\ 这样有用且强大的工具。这些作者帮助我开始了 \LaTeX；我希望为某个想开始使用 \ConTeXt\ 的人做同样的事情，尽管在本文档的西班牙语原版中，我只针对缺乏自己语言文档的西班牙语世界。我希望这份文档能满足这一期望，与此同时，其他人已慷慨地提出将其翻译成其他语言，于是有了这个英文版。谢谢。

\page[no]\blank

\rightaligned{Joaquín Ataz-López}\\
\rightaligned{2020年夏季}

% 如果没有 \page，本前言的最后一页将获得目录的页眉文本。
\page

\stoptitle

\stopcomponent


%%% Local Variables:
%%% mode: ConTeXt
%%% eval: (auto-fill-mode 1)
%%% TeX-master: "../introCTX_eng.tex"
%%% coding: utf-8-unix
%%% End:
%%% vim:set filetype=context tw=72 : %%%